\section{Tích hợp các module}\label{sec:6}

Bản chụp lúc START chỉ gồm \sig{DESC_PTR} và chỉ đi vào Main FSM. FSM dựng job từ descriptor, cấp
cấu hình của job đang tính cho array và VPU, cấp bank cho SRAM, và cấp địa chỉ, số hàng, số word,
stride cùng bank trong từng lệnh DMA.

\subsection{Kết nối giữa các module}\label{sec:6.1}

\begin{spectable}{|L{0.3}|L{0.3}|L{0.4}|}{Kết nối giữa các module}{tab:6-1}{\textbf{Từ} & \textbf{Tới} & \textbf{Tín hiệu}}
RF \sig{o_start_pulse}, \sig{o_soft_req}, \sig{o_cfg_desc_ptr} & FSM \sig{i_start_pulse}, \sig{i_soft_req}, \sig{i_cfg_desc_ptr} & Xung start, xung soft reset và \sig{DESC_PTR} \\\hline

FSM \sig{o_done_pulse}, \sig{o_err_pulse}, \sig{o_err_code}, \sig{o_perf_load}, \sig{o_perf_comp}, \sig{o_perf_store}, \sig{o_soft_busy}, \sig{o_err_desc_addr}, \sig{o_desc_done_cnt} & RF \sig{i_*} cùng tên & Trạng thái, mã lỗi, bộ đếm hiệu năng \\\hline

FSM \sig{o_soft_clr} & \sig{i_soft_clr} của RF, DMA, SRAM, Array, VPU & Xung xóa đồng bộ một chu kỳ; FSM không có \sig{i_soft_clr} \\\hline

FSM \sig{o_rd_start}, \sig{o_rd_dst}, \sig{o_rd_base}, \sig{o_rd_rows}, \sig{o_rd_rwords}, \sig{o_rd_stride}, \sig{o_rd_bank}, \sig{o_rd_abort} & DMA \sig{i_rd_*} cùng tên & Lệnh kênh đọc \\\hline

FSM \sig{o_wr_start}, \sig{o_wr_fmt}, \sig{o_wr_base}, \sig{o_wr_rows}, \sig{o_wr_rwords}, \sig{o_wr_stride}, \sig{o_wr_bank}, \sig{o_wr_abort} & DMA \sig{i_wr_*} cùng tên & Lệnh kênh ghi \\\hline

DMA \sig{o_rd_busy}, \sig{o_rd_done}, \sig{o_rd_error}, \sig{o_desc_valid}, \sig{o_desc_idx}, \sig{o_desc_data}, \sig{o_wr_busy}, \sig{o_wr_done}, \sig{o_wr_error} & FSM \sig{i_*} cùng tên & Trạng thái hai kênh và 16 word descriptor \\\hline

DMA \sig{o_spm_wr_valid}, \sig{o_spm_wr_buf}, \sig{o_spm_wr_bank}, \sig{o_spm_wr_addr}, \sig{o_spm_wr_data} & SRAM \sig{i_dw_valid}, \sig{i_dw_buf}, \sig{i_dw_bank}, \sig{i_dw_addr}, \sig{i_dw_data} & Word dữ liệu, địa chỉ, buffer, bank \\\hline

DMA \sig{o_spm_rd_valid}, \sig{o_spm_rd_bank}, \sig{o_spm_rd_addr} & SRAM \sig{i_dr_valid}, \sig{i_dr_bank}, \sig{i_dr_addr} & Yêu cầu đọc bank C \\\hline

SRAM \sig{o_dw_ready}, \sig{o_dr_ready}, \sig{o_dr_rvalid}, \sig{o_dr_rdata} & DMA \sig{i_spm_wr_ready}, \sig{i_spm_rd_ready}, \sig{i_spm_rd_rvalid}, \sig{i_spm_rd_rdata} & \sig{o_dw_ready} và \sig{o_dr_ready} là hằng 1; dữ liệu về sau đúng một chu kỳ \\\hline

FSM \sig{o_pe_start}, \sig{o_comp_dim} & Array \sig{i_pe_start}, \sig{i_matrix_dim} & Lệnh bắt đầu COMPUTE và \{8'h00, K, N, M\} \\\hline

Array \sig{o_pe_done}, \sig{o_busy} & FSM \sig{i_pe_done}, \sig{i_pe_busy} & \\\hline

FSM \sig{o_comp_post_en}, \sig{o_comp_relu_en}, \sig{o_comp_qm0}, \sig{o_comp_qshift} & VPU \sig{i_post_en}, \sig{i_relu_en}, \sig{i_quant_m0}, \sig{i_quant_shift} & Cấu hình hậu xử lý của job đang tính \\\hline

FSM \sig{o_comp_bank_a}, \sig{o_comp_bank_b}, \sig{o_comp_bank_bias}, \sig{o_comp_bank_c}, \sig{o_comp_active} & SRAM \sig{i_comp_bank_a}, \sig{i_comp_bank_b}, \sig{i_comp_bank_bias}, \sig{i_comp_bank_c}, \sig{i_comp_active} & Bank cấp cho phía COMPUTE \\\hline

Array \sig{o_a_rd_en}, \sig{o_a_rd_row}, \sig{o_a_rd_k}, \sig{o_b_rd_en}, \sig{o_b_rd_k}, \sig{o_b_rd_col} & SRAM \sig{i_a_rd_*}, \sig{i_b_rd_*}; về Array \sig{o_a_rd_data}, \sig{o_b_rd_data} & Đọc vector 64 bit; dữ liệu trả về trễ đúng một chu kỳ. \sig{o_a_rd_en} luôn bằng \sig{o_b_rd_en} \\\hline

Array \sig{o_row_valid}, \sig{o_row_idx}, \sig{o_row_col}, \sig{o_row_data}, \sig{o_row_mask} & VPU \sig{i_row_valid}, \sig{i_row_idx}, \sig{i_row_col}, \sig{i_row_data}, \sig{i_row_mask} & Luồng hàng 256 bit; \textbf{không} có đường ready \\\hline

VPU \sig{o_bias_rd_en}, \sig{o_bias_rd_col}, \sig{o_c_wr_en}, \sig{o_c_wr_fmt}, \sig{o_c_wr_row}, \sig{o_c_wr_col}, \sig{o_c_wr_data}, \sig{o_c_wr_mask} & SRAM \sig{i_bias_rd_*}, \sig{i_c_wr_*}; về VPU \sig{o_bias_rd_data} & Đọc bias 256 bit, ghi bank C 256 bit có mask \\\hline

top \sig{s_axi_*} & RF & 17 tín hiệu, đi thẳng, không logic chèn giữa \\\hline

DMA \sig{m_axi_*} & top & 25 tín hiệu, đi thẳng, không logic chèn giữa \\\hline

\end{spectable}

\subsection{Quy tắc tích hợp}\label{sec:6.2}

\paragraph{Tích hợp}\label{par:6-2}
\begin{itemize}
\item Tham số P phải thống nhất giữa Array, VPU, SRAM và trường P của \sig{HW_CONFIG}; baseline chốt P = 8 cho cả bốn chỗ
\item \sig{L_VPU} của Array phải bằng LAT của VPU, cả hai lấy từ \sig{gemm_pkg::VPU_LATENCY} = 4. Mỗi module kiểm điều kiện này lúc elaboration và phát \sig{\$error} nếu sai
\item Top truyền giá trị tham số đã chọn xuống các module, không phụ thuộc mặc định riêng của từng module
\item Top chỉ đấu nối và chứa bộ đồng bộ reset 2 flip-flop; không có logic dữ liệu, không có FSM, độ trễ xuyên qua top là 0 chu kỳ
\item \sig{rst_n_sync} là reset duy nhất cấp cho cả sáu module. \sig{soft_clr} đi tới năm module (RF, DMA, SRAM, Array, VPU) và không đi vào chân reset của bất kỳ flip-flop nào
\item Thiết kế có đúng ba nhóm đường tổ hợp từ input tới output của module: (1) VPU: \sig{o_bias_rd_en} = \sig{i_row_valid} \& \sig{i_post_en} và \sig{o_bias_rd_col} = \sig{i_row_col}; (2) DMA: \sig{m_axi_rready} = trạng thái RD\_DATA ∧ \sig{i_spm_wr_ready}, và vì \sig{o_dw_ready} là hằng 1 nên \sig{m_axi_rready} không phụ thuộc input AXI; (3) DMA: \sig{o_spm_wr_valid}, \sig{o_spm_wr_data} và (với FETCH) \sig{o_desc_valid}, \sig{o_desc_data} tổ hợp từ \sig{m_axi_rvalid}, \sig{m_axi_rdata} (đi vào SRAM hoặc bộ đệm descriptor của Main FSM, kết thúc ở thanh ghi/RAM). Ràng buộc I/O và kiểm tra vòng tổ hợp trong STA/lint phải bao phủ cả ba. Ready của Register File (\sig{s_axi_awready}, \sig{s_axi_wready}, \sig{s_axi_arready}) chỉ phụ thuộc thanh ghi nội bộ, không phụ thuộc tổ hợp vào input. Mọi output khác của mọi module đều đã qua thanh ghi
\item Không có đường tổ hợp từ input tới output trên cùng một cổng AXI
\end{itemize}

\subsection{Quy tắc giao tiếp}\label{sec:6.3}

\begin{spectable}{|L{0.24}|L{0.76}|}{Quy ước bắt tay giữa các module}{tab:6-3}{\textbf{Loại giao tiếp} & \textbf{Quy ước}}
Xung một chu kỳ & \sig{o_start_pulse}, \sig{o_soft_req}, \sig{o_soft_clr}, \sig{o_done_pulse}, \sig{o_err_pulse}, \sig{o_rd_done}, \sig{o_rd_error}, \sig{o_wr_done}, \sig{o_wr_error}, \sig{o_pe_done}, \sig{o_pe_start}: bằng 1 đúng một chu kỳ, bên nhận phải bắt ngay ở chu kỳ đó \\\hline

Lệnh start & \sig{o_rd_start} và \sig{o_wr_start} là xung một chu kỳ; payload lệnh (\sig{base}, \sig{rows}, \sig{rwords}, \sig{stride}, \sig{bank}, \sig{dst}/\sig{fmt}) phải hợp lệ \textbf{ở đúng chu kỳ đó}. Start khi kênh đang bận bị bỏ qua \\\hline

Mức (level) & \sig{o_rd_abort}, \sig{o_wr_abort}, \sig{o_comp_active}, \sig{o_comp_bank_*}, \sig{o_soft_busy}, \sig{i_rd_busy}, \sig{i_wr_busy}, \sig{i_pe_busy}: tín hiệu mức, giữ ổn định suốt khoảng có hiệu lực \\\hline

Cấu hình job & \sig{o_comp_dim}, \sig{o_comp_post_en}, \sig{o_comp_relu_en}, \sig{o_comp_qm0}, \sig{o_comp_qshift}: giữ ổn định từ trước \sig{o_pe_start} tới khi \sig{o_pe_done}. Assertion của VPU kiểm \sig{\$stable} trong khi đường ống còn bit valid \\\hline

Valid không ready & Luồng hàng Array → VPU (\sig{o_row_valid} …) và luồng descriptor DMA → FSM (\sig{o_desc_valid} …) không có đường ready: bên nhận \textbf{bắt buộc} nhận mỗi chu kỳ có valid. Điều này an toàn vì VPU có thông lượng 1 hàng/chu kỳ và bộ đệm descriptor của FSM luôn trống khi FETCH chạy \\\hline

Valid có ready & Đường DMA ↔ SRAM dùng valid/ready, nhưng \sig{o_dw_ready} và \sig{o_dr_ready} là hằng 1, nên không bao giờ có backpressure. Hiện thực vẫn phải giữ đúng ngữ nghĩa valid/ready \\\hline

Đọc SRAM & Mọi cổng đọc SRAM có độ trễ \textbf{đúng 1 chu kỳ}: dữ liệu hợp lệ ở chu kỳ ngay sau chu kỳ có \sig{*_rd_en}. Không có đường tổ hợp từ \sig{*_rd_en} tới dữ liệu \\\hline

\end{spectable}

\subsection{Xử lý lỗi giữa các module}\label{sec:6.4}

\begin{spectable}{|L{0.22}|L{0.18}|L{0.6}|}{Nguồn lỗi và đường báo lỗi}{tab:6-4}{\textbf{Khối} & \textbf{Có lỗi?} & \textbf{Cách báo}}
Register File & Không & Mọi truy cập AXI4-Lite trả OKAY, kể cả offset dự trữ. RF chỉ \emph{lưu} \sig{ERR_CODE} do FSM cấp \\\hline

Main FSM & Có & Nguồn duy nhất của \sig{ERR_CODE}. Phát mã 1 từ kết quả kiểm tra descriptor; chuyển \sig{o_rd_error} thành mã 2 hoặc 4 tùy pha, \sig{o_wr_error} thành mã 3 \\\hline

DMA & Có & \sig{o_rd_error} và \sig{o_wr_error}, mỗi cái một xung. Không mang mã lỗi: FSM quyết mã theo pha đang chạy \\\hline

Local SRAM & Không & Không có output lỗi. Vi phạm địa chỉ và vi phạm bank chỉ bị bắt bằng assertion mô phỏng \\\hline

Systolic Array & Không & Không có output lỗi. Giả định M, N, K đã hợp lệ vì FSM kiểm trước khi phát \sig{o_pe_start} \\\hline

VPU & Không & Không có output lỗi. Tràn không xảy ra trên miền hợp lệ; bước tạo Y saturate thay vì báo lỗi \\\hline

\end{spectable}

\paragraph{Lan truyền lỗi}\label{par:6-5}
\begin{itemize}
\item DMA của kênh lỗi ngừng phát burst mới, nhận hết beat R hoặc phản hồi B của burst đã cam kết, rồi phát một xung lỗi. Beat dữ liệu của burst lỗi vẫn được ghi vào scratchpad; nội dung đó không được dùng
\item FSM ghi nhận lỗi vào một ô duy nhất gồm \sig{err_v}, \sig{err_idx}, \sig{err_code}, \sig{err_addr}; khi có lỗi mới, chỉ ghi đè nếu \sig{err_idx} nhỏ hơn giá trị đang giữ. Nhờ vậy \sig{ERR_CODE} và \sig{ERR_DESC_ADDR} luôn ứng với descriptor lỗi sớm nhất theo thứ tự chuỗi
\item Gọi f là descriptor lỗi sớm nhất. Sau khi có lỗi, FSM không bắt đầu pha nào của descriptor ≥ f (không lệnh DMA, không \sig{o_pe_start} cho chúng); các descriptor đứng trước f vẫn chạy tới hết, kể cả STORE\_C. Pha đang chạy chạy tới hết. Khi hệ yên — engine Đọc đã dừng và không chờ DMA, engine Tính không ở C\_RUN, engine Ghi rảnh và không còn job nào có chỉ số < f chờ tính hoặc ghi — FSM phát \sig{o_err_pulse}. Điều kiện chính xác nằm ở mục Kết thúc tác vụ của tài liệu Main FSM; nhìn từ bên ngoài, lúc đó \sig{i_rd_busy} = \sig{i_wr_busy} = \sig{i_pe_busy} = 0
\item Array và VPU không bị thông báo về lỗi; chúng chỉ ngừng vì không nhận thêm \sig{o_pe_start}. Nếu lỗi xảy ra trong lúc COMPUTE đang chạy, COMPUTE vẫn chạy tới \sig{o_pe_done} rồi hệ thống mới yên
\item Lỗi không bao giờ làm treo: mọi pha đang chạy đều có điều kiện kết thúc hữu hạn (burst AXI kết thúc bằng \sig{rlast} hoặc \sig{bvalid}; COMPUTE kết thúc sau số chu kỳ xác định)
\end{itemize}

\subsection{Kiểm chứng ở mức hệ thống}\label{sec:6.5}

\begin{spectable}{|L{0.12}|L{0.88}|}{Điểm kiểm chứng mức hệ thống}{tab:6-6}{\textbf{ID} & \textbf{Nội dung}}
SYS-01 & Nội dung DRAM sau một tác vụ khớp từng bit với mô hình tham chiếu ở mục~\ref{sec:3.5}, áp dụng lần lượt cho từng descriptor theo thứ tự chuỗi \\\hline

SYS-02 & Port top-level đúng 44 tín hiệu, đúng tên, độ rộng và chiều như bảng~\ref{tab:1-6} và bảng~\ref{tab:1-7}; không có port nào khác \\\hline

SYS-03 & Mỗi descriptor được FETCH đúng một lần, bằng một burst 16 beat, theo đúng thứ tự chuỗi \\\hline

SYS-04 & Pha STORE\_C ra DRAM đúng theo thứ tự descriptor; \sig{DESC_DONE_CNT} tăng đơn điệu từng đơn vị \\\hline

SYS-05 & Tại mọi chu kỳ, mỗi bank của mỗi buffer có tối đa một bên truy cập (assertion AS-02 của SRAM) \\\hline

SYS-06 & Sau khi \sig{o_pe_done}, đường ống VPU không còn bit valid nào (assertion AS-07 của VPU) \\\hline

SYS-07 & \sig{o_done_pulse} và \sig{o_err_pulse} không bao giờ cùng một chu kỳ; sau khi START được chấp nhận, đúng một trong ba bit BUSY, DONE, ERROR bằng 1 ở mọi lần đọc \sig{STATUS} \\\hline

SYS-08 & Khi lỗi, \sig{ERR_DESC_ADDR} và \sig{DESC_DONE_CNT} ứng với descriptor lỗi sớm nhất; không pha nào của descriptor ≥ f đã bắt đầu \\\hline

SYS-09 & Chuỗi soft reset luôn kết thúc: từ \sig{o_soft_req} tới \sig{o_soft_clr} là hữu hạn chu kỳ, và sau đó \sig{STATUS} = 0 mà không có DONE và không có ERROR \\\hline

SYS-10 & Kết quả không phụ thuộc mức độ chạy chồng: chạy cùng chuỗi với bus chậm (nhiều bubble AXI) và bus nhanh cho cùng nội dung DRAM \\\hline

SYS-11 & \sig{DESC_PTR} không căn 64 byte kết thúc tác vụ với \sig{ERR_CODE} = 1 mà \textbf{không} có giao dịch AXI master nào và không có truy cập scratchpad nào \\\hline

SYS-12 & Mọi truy cập AXI4-Lite, kể cả vào offset dự trữ 0x08 … 0x24, trả \sig{s_axi_bresp} / \sig{s_axi_rresp} = 2'b00 \\\hline

\end{spectable}

\begin{notebox}[Phạm vi]
Bảng~\ref{tab:6-6} là \textbf{yêu cầu kiểm chứng} của tài liệu thiết kế, không phải báo cáo kết quả.
Bộ tài liệu baseline này là tài liệu; nó chưa được đối chiếu với RTL và chưa có mô phỏng nào được
chạy.
\end{notebox}
